\chapter{Analyse des besoins et cahier des charges}
\label{ch:besoins}

\section{Présentation du contexte}

Le cahier des charges se compose de 21 sections décrivant, dans l'ordre, les objectifs du
système, la typologie des utilisateurs, la gestion des profils, les données temporelles,
la gouvernance, puis quinze modules fonctionnels, et se termine par un ensemble de règles
générales. Sa lecture fait apparaître une structure en trois couches : une couche
d'identité et de gouvernance, une couche d'activité scientifique, et une couche de
ressources et de traçabilité. Ce découpage a été conservé jusque dans l'organisation des
packages Java et dans la répartition des diagrammes de classes.

Une caractéristique de l'énoncé mérite d'être relevée : il ne décrit jamais d'écran. Il
décrit des règles de fonctionnement institutionnel --- « historique complet conservé »,
« uniquement par le directeur en mandat », « qui a reçu / qui n'a pas reçu ». L'analyse a
donc consisté à traduire des règles de gestion en contraintes vérifiables par le système,
et non à transposer une maquette.

\section{Acteurs du système}

Le cahier des charges distingue les visiteurs des membres, et détaille cinq catégories de
membres ainsi que trois responsabilités de gouvernance et un administrateur technique. Le
tableau~\ref{tab:acteurs} recense les acteurs effectivement implémentés, avec le rôle
applicatif correspondant dans l'énumération \code{RoleName}.

\begin{table}[H]
\centering
\caption{Acteurs du système et rôles applicatifs correspondants}
\label{tab:acteurs}
\small
\begin{tabularx}{\textwidth}{@{}p{3.1cm}p{3.5cm}X@{}}
\toprule
\textbf{Acteur} & \textbf{\code{RoleName}} & \textbf{Périmètre d'accès} \\
\midrule
Visiteur & --- (anonyme) & Portail public, dépôt d'une candidature d'adhésion \\
Doctorant & \code{PHD\_STUDENT} & Profil, membres, publications (ajout et consultation), événements, calendrier, notifications. \emph{Aucun accès aux modules internes} \\
Membre associé & \code{ASSOCIATE\_MEMBER} & Le précédent, plus documents, matériel et procès-verbaux \\
Membre permanent & \code{PERMANENT\_MEMBER} & Identique au membre associé \\
Chef d'équipe & \code{TEAM\_LEADER} & Le précédent, plus création d'événements et de réunions, dépôt de procès-verbaux \\
Vice-directeur & \code{VICE\_DIRECTOR} & Le précédent, plus validation des demandes de matériel, statuts, conventions, rapports, journal d'audit, inventaire \\
Directeur & \code{DIRECTOR} & Le précédent, plus décision sur les adhésions, gestion des équipes et attribution des mandats \\
Retraité & \code{RETIRED\_MEMBER} & Compte gelé : authentification refusée, historique conservé \\
Ancien membre & \code{FORMER\_MEMBER} & Compte désactivé, réactivable \\
Administrateur système & \code{SYSTEM\_ADMIN} & Administration technique : comptes, rôles, identité du laboratoire. Pas de compétence métier \\
\bottomrule
\end{tabularx}
\end{table}

La distinction entre l'administrateur système et le directeur mérite d'être soulignée : le
cahier des charges leur attribue des responsabilités de nature différente (§5.1 et §5.2).
L'administrateur gère la technique, le directeur gouverne le laboratoire. Cette séparation
a été conservée jusque dans les annotations de sécurité : l'administrateur système ne peut
pas statuer sur une demande d'adhésion, et le directeur ne peut pas créer de compte
technique.

La figure~\ref{fig:usecase} présente le diagramme de cas d'utilisation. Il a été construit
à partir des annotations \code{@PreAuthorize} réellement présentes sur les contrôleurs et
du fichier de garde \code{utils/roles.js} côté client : chaque flèche correspond donc à un
endpoint effectivement exposé et autorisé, et non à une intention.

\begin{figure}[H]
\centering
\includegraphics[width=\textwidth,height=0.80\textheight,keepaspectratio]{\figpath/uml/01-use-case.png}
\caption{Diagramme de cas d'utilisation du système LIAS}
\label{fig:usecase}
\end{figure}

L'organisation du diagramme traduit une relation d'héritage entre acteurs : le membre
hérite du visiteur, le chef d'équipe du membre, le vice-directeur du chef d'équipe et le
directeur du vice-directeur. Chaque niveau ajoute des cas d'utilisation sans en retirer.
L'administrateur système est délibérément placé hors de cette chaîne, avec ses deux seuls
cas d'utilisation propres.

\section{Besoins fonctionnels}

Les besoins fonctionnels sont issus des sections~3 à~20 du cahier des charges. Le
tableau~\ref{tab:besoinsfonc} les regroupe en huit domaines et précise l'exigence
caractéristique de chacun --- celle qui, si elle n'était pas satisfaite, rendrait le
module non conforme.

\begin{table}[H]
\centering
\caption{Besoins fonctionnels par domaine et exigence caractéristique}
\label{tab:besoinsfonc}
\small
\begin{tabularx}{\textwidth}{@{}p{1.1cm}p{3.5cm}X@{}}
\toprule
\textbf{§ CdC} & \textbf{Domaine} & \textbf{Exigence caractéristique} \\
\midrule
3, 4    & Profils et données temporelles & Un membre peut quitter le laboratoire et revenir ; chaque période d'affiliation est conservée séparément \\
5, 7    & Gouvernance et rôles & Les mandats définissent des périodes datées ; l'historique complet des rôles est conservé \\
6       & Demandes d'adhésion & La décision revient \emph{uniquement} au directeur en mandat ; l'acceptation crée le compte \\
8       & Départ et retraite & Le compte est gelé ou désactivé, jamais supprimé ; la réactivation reste possible \\
9, 12, 17, 20 & Activité scientifique & Les événements sont archivés \emph{par édition} ; les publications sont classables par auteur, année et équipe \\
10      & Archivage documentaire & Les documents sont structurés par événement et par date, et historisés en versions \\
11      & Rapport annuel & Génération \emph{automatique} à partir des événements, publications et activités \\
13--16, 18, 19 & Ressources et communication & Le suivi du matériel doit indiquer qui a reçu \emph{et qui n'a pas reçu} \\
\bottomrule
\end{tabularx}
\end{table}

Deux exigences ont eu un effet structurant disproportionné par rapport à leur longueur
dans l'énoncé. La première est l'archivage par édition (§9) : elle impose que deux
occurrences successives d'un même colloque coexistent sans écrasement, ce qui se traduit
par un attribut \code{editionNumber} et un endpoint dédié. La seconde est le suivi des
non-bénéficiaires de matériel (§14) : elle exige une requête calculant un complément
d'ensemble --- les membres actifs n'ayant pas reçu un article donné --- et non une simple
lecture de table.

\section{Besoins non fonctionnels}

Le document de spécification associé au cahier des charges énumère un ensemble d'exigences
non fonctionnelles. Le tableau~\ref{tab:nonfonc} en donne l'état réel, sans complaisance :
une exigence non satisfaite y est déclarée comme telle et reprise dans les limites
(§\ref{sec:limites}).

\begin{table}[H]
\centering
\caption{Besoins non fonctionnels et niveau de satisfaction atteint}
\label{tab:nonfonc}
\small
\begin{tabularx}{\textwidth}{@{}p{3.9cm}p{1.9cm}X@{}}
\toprule
\textbf{Exigence} & \textbf{État} & \textbf{Élément justificatif} \\
\midrule
Sécurité & Satisfait & JWT signé HS512, BCrypt (coût 12), contrôle d'accès par rôle appliqué méthode par méthode, gel des comptes vérifié à chaque requête \\
Architecture propre & Satisfait & 18 packages métier, séparation contrôleur~/~service~/~repository~/~DTO systématique \\
API REST & Satisfait & 22 contrôleurs sous \code{/api/v1}, documentation OpenAPI générée \\
Composants réutilisables & Satisfait & Système de composants interne (8 modules) réutilisé par les 21 pages \\
Design responsive & Satisfait & Navigation en tiroir sous 1024~px, tables à défilement horizontal encapsulé \\
Interface professionnelle & Satisfait & Identité visuelle unique appliquée à l'ensemble des écrans \\
Maintenabilité & Satisfait & Découpage constant, 91 tests verrouillant les règles critiques \\
Performance & Partiel & \code{open-in-view} désactivé, pagination serveur, index SQL ; aucune campagne de charge menée \\
Accessibilité & Partiel & Libellés accessibles sur les boutons-icônes, focus visible, fermeture clavier des modales ; pas d'audit WCAG formel \\
Thème sombre & \textbf{Non satisfait} & Non implémenté --- limite assumée (§\ref{sec:limites}) \\
\bottomrule
\end{tabularx}
\end{table}

\section{Contraintes}

Quatre contraintes ont encadré la réalisation.

\paragraph{Contrainte académique.} Le projet relève d'un module Java Enterprise Edition :
la plateforme technique du backend est imposée dans son principe. Le choix de Spring~Boot
comme implémentation des spécifications JEE est justifié au chapitre~\ref{ch:techno}.

\paragraph{Contrainte de conservation.} La règle « aucun membre ne perd ses données »
interdit la suppression physique. Elle se traduit par une suppression logique généralisée
au niveau de la classe de base des entités, et impose de traiter partout le cas d'une
entité archivée encore référencée par une autre.

\paragraph{Contrainte d'authenticité des données.} Le document de spécification proscrit
explicitement les maquettes, les données fictives et les pages de remplissage. Chaque
écran présenté dans ce rapport est connecté à des endpoints réels interrogeant la base de
données.

\paragraph{Contrainte d'exécution.} L'application devait fonctionner sur une base
PostgreSQL hébergée, avec un schéma versionné et validé au démarrage --- ce qui interdit
la génération automatique du schéma par l'ORM et impose des migrations explicites.

\section{Critères de réussite}

Les critères retenus pour juger le projet abouti, et vérifiés au chapitre~\ref{ch:tests},
sont les suivants :

\begin{enumerate}
  \item les 21 sections du cahier des charges sont couvertes par une fonctionnalité
        opérationnelle, et non par une page vide ;
  \item la séparation \srt{}~\neqq{}~\rol{}~\neqq{}~\man{} est vérifiable : aucun compte
        ne présente d'écart entre son rôle stocké et le rôle que son statut et ses mandats
        impliquent ;
  \item les permissions du client et celles du serveur coïncident : aucune commande
        visible ne doit aboutir à un refus, aucune opération autorisée ne doit être
        masquée ;
  \item un compte retraité ou ancien ne peut plus accéder au système, alors que son
        historique demeure consultable ;
  \item la compilation, la construction du client et la suite de tests s'exécutent sans
        erreur.
\end{enumerate}
